% !TeX root = ../../Book4.tex
% 纲要标题：### 20.1 标准截断
% 纲要标签：b4:sec:1901
% 纲要位置：计划纲要.md，第 3742 行。

\section{标准截断}
\label{b4:sec:1901}

截去复形的一端会改变截口处的核或像。典范截断修正端点项，既保留指定范围的上同调，又将原对象放入一个截断三角。


% 纲要内容（仅作写作计划，尚非教材正文）：
%
% * 朴素截断与典范截断
% * 标准 t-结构
% * 截断三角
%

\subsection{朴素截断与典范截断}
\label{b4:subsec:1901:01}

把复形截成两段，首先遇到的问题是截口的微分。直接删掉高次项会扩大最后一项的圈；直接删掉低次项则会丢掉原来进入第一项的边界。这两种操作适合描述复形的项，却不能原样保留指定次数的上同调。

\begin{definition}\label{b4:def:standard-truncations}
设 \(C\) 为交换范畴 \(\mathcal A\) 中的上链复形，\(n\in\mathbb Z\)。保留次数 \(i\le n\) 的原项与原微分、把其余项置零，所得复形记为 \(\sigma^{\le n}C\)；相应地定义 \(\sigma^{\ge n}C\)。这些构造称为\term[pusujianduan]{朴素截断}[brutal truncation]。把截口改为
\[
 (\tau^{\le n}C)^i=
 \begin{cases}C^i&i<n,\\ \ker \mathrm{d}^n&i=n,\\0&i>n,\end{cases}
 \qquad
 (\tau^{\ge n}C)^i=
 \begin{cases}0&i<n,\\ \operatorname{coker}\mathrm{d}^{n-1}&i=n,\\C^i&i>n,\end{cases}
\]
并采用原微分诱导的映射，得到\term[dianfanjianduan]{典范截断}[canonical truncation]，也称好截断。它们分别带有自然链映射 \(\tau^{\le n}C\to C\) 和 \(C\to\tau^{\ge n}C\)。
\symidx{\(\tau^{\le n},\tau^{\ge n}\)：典范截断函子}
\end{definition}
\symidx{\(\sigma^{\le n},\sigma^{\ge n}\)：上链复形的朴素截断}

复形映射在截口的核上限制、在余核上诱导映射，因而两种典范截断都是复形范畴上的函子。它们也保持链同伦。具体地，设 \(f,g:C\to D\) 及次数为 \(-1\) 的 \(h\) 满足
\[
f^i-g^i=\mathrm d_D^{i-1}h^i+h^{i+1}\mathrm d_C^i.
\]
记 \(\iota_C^n:Z^n(C)\hookrightarrow C^n\) 为核包含。对 \(\tau^{\le n}\)，低于截口的同伦仍取 \(h^i\)，在截口取
\(\bar h^n=h^n\iota_C^n:Z^n(C)\to D^{n-1}\)，其余取零。因为
\(\mathrm d_C^n\iota_C^n=0\)，原同伦等式在截口变为
\[
(f^n-g^n)\iota_C^n
=\mathrm d_D^{n-1}h^n\iota_C^n.
\]
两边都经 \(Z^n(D)\) 因子化，核包含的单射性便给出截断后的同伦等式；在 \(n-1\) 次，微分经 \(Z^n(C)\) 的因子化也使原等式保持。

\ProseParagraphBreak
对 \(\tau^{\ge n}\)，记 \(\pi_C^n:C^n\twoheadrightarrow\operatorname{coker}\mathrm d_C^{n-1}\)、\(\pi_D^n:D^n\twoheadrightarrow\operatorname{coker}\mathrm d_D^{n-1}\) 为商映射。同伦在 \(i\le n\) 时取零，在 \(n+1\) 次取
\(\bar h^{n+1}=\pi_D^n h^{n+1}\)，在更高次数仍取 \(h^i\)。若 \(\bar f^n,\bar g^n,\bar{\mathrm d}_C^n\) 表示截口的诱导映射，则
\[
\begin{aligned}
(\bar f^n-\bar g^n)\pi_C^n
&=\pi_D^n(f^n-g^n)\\
&=\pi_D^n h^{n+1}\mathrm d_C^n
=\bar h^{n+1}\bar{\mathrm d}_C^n\pi_C^n.
\end{aligned}
\]
这里 \(\pi_D^n\mathrm d_D^{n-1}=0\) 消去了另一项；\(\pi_C^n\) 为满态射，所以截口的同伦等式成立。在 \(n+1\) 次，恒等式
\(\bar{\mathrm d}_D^n\pi_D^n=\mathrm d_D^n\) 恢复原来的同伦等式，其他次数不变。因此两种截断都定义了 \(K(\mathcal A)\) 上的函子。

\ProseParagraphBreak
上截断在截口的核是整个 \(\ker\mathrm{d}_C^n\)，进入它的像仍为 \(\operatorname{im}\mathrm{d}_C^{n-1}\)，所以
\[
H^n(\tau^{\le n}C)
=\frac{\ker\mathrm{d}_C^n}{\operatorname{im}\mathrm{d}_C^{n-1}}
=H^n(C).
\]
下截断在截口使用 \(C^n/\operatorname{im}\mathrm{d}_C^{n-1}\)。其诱导微分
\(\bar{\mathrm{d}}^n:C^n/\operatorname{im}\mathrm{d}_C^{n-1}\rightarrow C^{n+1}\)
的核为 \(\ker\mathrm{d}_C^n/\operatorname{im}\mathrm{d}_C^{n-1}\)，而截口前一项为零，故
\[
H^n(\tau^{\ge n}C)=\ker\bar{\mathrm{d}}^n\cong H^n(C).
\]
其他保留次数的核与像不变，被删去次数的上同调为零。因此上截断保留 \(i\le n\) 的上同调，下截断保留 \(i\ge n\) 的上同调；核与余核的端点修正正是为了保留截口的原上同调。它们由此保持拟同构，将上述同伦范畴上的函子与局部化函子复合，依\cref{b4:def:derived-category}的泛性质得到导出范畴上的截断函子。

\ProseParagraphBreak
例如 \(C=(\mathbb Z\xrightarrow{2}\mathbb Z)\) 位于次数 \(-1,0\)。朴素截断 \(\sigma^{\ge0}C=\mathbb Z[0]\)，而典范截断 \(\tau^{\ge0}C=(\mathbb Z/2\mathbb Z)[0]\)。前者忘记了进入零次项的边界，后者才保留 \(H^0(C)\)。

\subsection{标准 t-结构}
\label{b4:subsec:1901:02}

\begin{definition}\label{b4:def:standard-t-halves}
设 \(\mathcal A\) 为交换范畴，\(D(\mathcal A)\) 为其导出范畴，\(n\in\mathbb Z\)。令
\[
 D^{\le n}(\mathcal A)=\{X:H^i(X)=0\ (i>n)\},\qquad
 D^{\ge n}(\mathcal A)=\{X:H^i(X)=0\ (i<n)\}
\]
表示由相应对象组成的全子范畴。它们也分别由同构于次数不超过 \(n\)、次数不低于 \(n\) 的复形的对象组成。\((D^{\le0},D^{\ge0})\) 称为\term[biaozhuntjiegou]{标准 t-结构}[standard t-structure]；下一节将抽取它满足的公理。
\symidx{\(D^{\le n},D^{\ge n}\)：标准上同调次数条件}
\end{definition}

\begin{proposition}\label{b4:prop:standard-truncation-orthogonality}
设 \(\mathcal A\) 为局部小交换范畴，\(n\in\mathbb Z\)。若 \(X\in D^{\le n}(\mathcal A)\)、\(Y\in D^{\ge n+1}(\mathcal A)\)，则 \(\operatorname{Hom}_{D(\mathcal A)}(X,Y)=0\)。
\end{proposition}
\begin{proof}
以典范截断把 \(X,Y\) 分别换成项集中在 \(i\le n\)、\(i\ge n+1\) 的复形。这只能说明直接的链映射 \(X\to Y\) 为零；导出态射还可能经由屋顶的中间复形，而它未必有相同的项次数界。用屋顶 \(X\xleftarrow{s}Z\xrightarrow{f}Y\) 表示任意导出态射。因为 \(s\) 为拟同构且 \(X\) 的高次上同调为零，\(\tau^{\le n}Z\to Z\) 也是拟同构。用它细化得到
\[
X\xleftarrow{\ \simeq\ }\tau^{\le n}Z\longrightarrow Y.
\]
此时中间复形的项才集中在 \(i\le n\)，向 \(Y\) 的链映射逐项为零：源、靶没有共同的非零次数。因此每个屋顶都表示零态射，正交性成立。
\end{proof}

这个零性要求源在低次、靶在高次。反向一般不为零：把零次对象 \(M\) 移位为 \(M[1]\)，它位于 \(-1\) 次，而 \(\operatorname{Hom}(N,M[1])\) 可以记录一次扩张。这里的次数判断始终使用 \(C[1]^i=C^{i+1}\)。

\subsection{截断三角}
\label{b4:subsec:1901:03}

\begin{proposition}\label{b4:prop:standard-truncation-triangle}
设 \(C\) 为交换范畴 \(\mathcal A\) 中的复形。对每个整数 \(n\)，导出范畴中有自然的好三角
\[
 \tau^{\le n}C\longrightarrow C\longrightarrow
 \tau^{\ge n+1}C\longrightarrow(\tau^{\le n}C)[1].
\]
两端分别属于 \(D^{\le n}\) 和 \(D^{\ge n+1}\)。
\end{proposition}
\begin{proof}
取逐项短正合列 \(0\to\tau^{\le n}C\to C\to Q\to0\)。商复形 \(Q\) 在 \(n\) 次为 \(\operatorname{im}\mathrm{d}^n\)，在较高次数与 \(C\) 相同，较低次数为零。从 \(Q\) 到 \(\tau^{\ge n+1}C\) 的映射在 \(n\) 次为零，在 \(n+1\) 次取商，在更高次数为恒等。它的核是两项恒等复形 \(\operatorname{im}\mathrm{d}^n\to\operatorname{im}\mathrm{d}^n\)，所以是拟同构。由\cref{b4:prop:derived-short-exact-triangle}得到的三角，把 \(Q\) 换成这一同构对象即得结论。所有映射来自核、商及原链映射，故对 \(C\) 自然。
\end{proof}

三角的第三箭头不应省去。若 \(C\) 只在 \(n,n+1\) 两次有上同调，该箭头仍可能阻止它分裂为两个上同调对象的直和。截断提供逐层分析的方法，并未断言上同调完全决定复形的导出同构类型。标准截断与下面的抽象构造可对照 Beilinson、Bernstein、Deligne 的《Faisceaux pervers》\cite[Exemples 1.3.2; Proposition 1.3.3]{BBD1982FaisceauxPervers}。
